% ============================================================
%  HARDWARE SECTION  --  GymEMG-Net
%  Figures required (add to figures/ directory):
%    fig:arch            -- system_architecture.png
%    fig:schematic       -- node_schematic.png
%    fig:prototype       -- node_prototype.jpg
%    fig:enclosure       -- node_enclosure.jpg
%    fig:placement       -- electrode_placement_overview.jpg
%    fig:placement_zoom  -- electrode_placement_detail.jpg
%    fig:pipeline        -- signal_pipeline.png
% ============================================================

\section{Hardware Design}

\subsection{System Architecture}

Centralised, wired sEMG acquisition introduces a set of problems that become increasingly difficult to manage in dynamic exercise environments: electrode lead wires pick up motion artefact proportional to cable length, restrict the subject's range of motion, and impose a physical tether that is impractical outside a laboratory. To avoid these trade-offs, the acquisition electronics in GymEMG-Net are distributed directly to the muscles.

The system comprises four identical battery-powered sensor nodes and one master gateway arranged in a star topology, as shown in Fig.~\ref{fig:arch}. Each node is positioned over a target muscle group and is entirely self-contained---it handles local amplification, digitisation, and wireless transmission without any wired connection to the host. The four slave nodes communicate with the master ESP32 over ESP-NOW~\cite{espressif_espnow}; the master aggregates all four channels and relays the merged stream to a host PC via USB-UART, where a monitoring GUI renders the live signals in real time and logs them to disk for offline analysis. No access point, router, or Bluetooth pairing is involved at any stage of the link.

\begin{figure}[!t]
\centering
\includegraphics[width=\columnwidth]{figures/system_architecture.png}
\caption{Overall system architecture. Four wireless sensor nodes (each comprising a MyoWare~2.0 and an ESP32-WROOM-32) transmit raw sEMG frames to a master ESP32 over ESP-NOW. The master aggregates all four channels and forwards the merged stream to a host PC via USB serial, where a monitoring GUI renders the live signals and exports session data.}
\label{fig:arch}
\end{figure}


\subsection{Sensor Node Design}

\subsubsection{sEMG Front-End: MyoWare 2.0}

The MyoWare 2.0 (Advancer Technologies) integrates a differential instrumentation amplifier, a band-pass filter stage, and a full-wave rectifier onto a compact 44.7\,mm\,$\times$\,20.3\,mm board~\cite{advancer_myoware2}. The key design advantage for a body-worn application is its snap-electrode interface: three integrated connectors accept standard Ag/AgCl disposable pads directly on the board, eliminating the intermediate lead wires that are a primary source of motion-artefact contamination in ambulatory recordings. The two in-line snaps form the differential recording pair, positioned approximately 20\,mm apart along the muscle belly in accordance with SENIAM guidelines~\cite{Hermens2000SENIAM}. The third connector (REF) accepts a passive reference electrode placed over a bony landmark and is connected to the board's analogue ground, providing common-mode rejection without requiring a driven-right-leg circuit.

The RAW output pin, which carries the amplified, band-pass-filtered bipolar sEMG signal prior to rectification, is routed directly to GPIO34 (ADC1\_CH6) of the ESP32. Using the pre-rectification output preserves the full spectral content of the signal---nominally 20--500\,Hz, as defined by the MyoWare~2.0 hardware band-pass filter---and gives the downstream classifier access to both amplitude and frequency features, rather than constraining it to the hardware-defined envelope time constant~\cite{phinyomark2018feature}. This approach is consistent with comparable low-cost wearable sEMG implementations~\cite{gomez2022wearable}. Recent work confirms that 1D-CNN classifiers operating on raw lower-limb sEMG achieve $97.0 \pm 0.8$\% accuracy across seven motion classes~\cite{mot_semg2025}, validating the spectral richness of the pre-rectification signal for gym exercise recognition.

% [AUTHOR NOTE: Consider including the gain setting or gain range of the
% MyoWare 2.0 instrumentation amplifier used in this configuration.]

\subsubsection{Microcontroller: ESP32 DevKit (WROOM-32)}

Each sensor node is built around an ESP32 DevKit development board (featuring the ESP32-WROOM-32 module)~\cite{espressif_esp32}, powered by a dual-core Xtensa LX6 microcontroller clocked at 240\,MHz with an on-chip 12-bit successive-approximation register (SAR) ADC. The MyoWare RAW output is sampled on ADC1\_CH6 (GPIO34) using \texttt{analogReadMilliVolts()}, which applies the ESP32 factory calibration to return a millivolt value directly.

The sampling rate of 1\,kSa/s is selected in accordance with the Nyquist--Shannon sampling theorem~\cite{shannon1949communication}, which requires the sampling frequency $f_s$ to satisfy
\begin{equation}
  f_s \geq 2 \cdot f_{\max},
\end{equation}
where $f_{\max}$ is the highest frequency of interest in the signal. The MyoWare~2.0 hardware band-pass filter passes nominally 20--500\,Hz; therefore the minimum permissible sampling rate is $2 \times 500 = 1000\,\text{Sa/s}$. Choosing $f_s = 1\,\text{kSa/s}$ exactly meets this Nyquist limit, and the 450\,Hz software low-pass filter applied downstream provides a 50\,Hz guard band (one decade of roll-off headroom) that prevents any spectral leakage from the 450--500\,Hz transition region from corrupting the signal band of interest. The effective number of bits (ENOB) at this operating point is approximately 11, consistent with published characterisation of the ESP32 ADC1 under similar supply and input conditions~\cite{gomez2022wearable}.

Samples are accumulated into fixed-length 50-point windows (50\,ms). Each batch packet is structured as follows: a 1-byte type identifier, a 1-byte node ID, a 4-byte frame counter (\texttt{frame\_id\_start}), a 4-byte frame-based millisecond timestamp (\texttt{t0\_ms}), a 4-byte inter-sample interval (\texttt{dt\_us}\,=\,1000\,$\mu$s), a 2-byte sample count, and 100\,bytes of unsigned 16-bit ADC millivolt values---approximately 116\,bytes in total, well within the 250-byte ESP-NOW payload ceiling. Each node transmits twenty batch packets per second ($1000\,\text{Sa/s} \div 50 = 20\,\text{batches/s}$), yielding a per-channel data throughput of $20 \times 116 \approx 2.3\,\text{kB/s}$.

A 3\,mm green LED, connected directly to the 3.3\,V rail through a 470\,$\Omega$ current-limiting resistor (R1), serves as a simple power-on indicator. It illuminates whenever the regulated supply is active, allowing the operator to confirm that the node is powered without requiring any software interaction.

\subsubsection{Power Conditioning: 3.3\,V Regulator and Decoupling}

The Li-Po cell output (nominally 3.7\,V, ranging from 3.0\,V to 4.2\,V over the discharge cycle) is regulated to a stable 3.3\,V by a low-dropout linear regulator before reaching the ESP32 and the MyoWare front-end. This ensures that ADC reference voltage variation does not introduce a supply-dependent offset into the sEMG readings. Two decoupling capacitors are placed at the regulator output: a 10\,$\mu$F bulk capacitor (C1) that suppresses low-frequency load transients caused by the ESP32 radio bursts, and a 100\,nF ceramic capacitor (C2) that attenuates high-frequency switching noise and provides local charge storage for the ADC sample-and-hold circuit. Together, C1 and C2 maintain a clean, stable supply to both the analogue front-end and the digital core throughout the wireless transmission cycle.

\subsubsection{Charging and Power Path}

The Li-Po cell is charged via a TP4056-based module (U1) with an on-board 4.2\,k$\Omega$ programming resistor (R2) that sets the constant-current charge rate to approximately 250\,mA. A USB-C receptacle on the module provides the charge input, allowing standard USB power banks or wall adapters to be used in the field. An SPST switch (S1) is positioned on the regulated 3.3\,V rail between the output of the LDO regulator and the 3.3\,V power input of the ESP32 DevKit and MyoWare front-end. This configuration allows the active components to be powered off while keeping the battery connected to the TP4056 charger module. The complete circuit schematic for one sensor node is shown in Fig.~\ref{fig:schematic}.

\begin{figure}[!t]
\centering
\includegraphics[width=\columnwidth]{figures/node_schematic.png}
\caption{Circuit schematic for a single GymEMG-Net sensor node. The TP4056 module (U1) charges the 3.7\,V Li-Po cell via USB-C; the regulated 3.3\,V output (via the LDO, with C1\,=\,10\,$\mu$F and C2\,=\,100\,nF decoupling) is routed through the SPST power switch (S1) to supply the ESP32 DevKit and the MyoWare~2.0 front-end. The MyoWare RAW output connects to GPIO34 (ADC1). The green power-indicator LED is current-limited by R1\,=\,470\,$\Omega$ and tied directly to the switched 3.3\,V rail.}
\label{fig:schematic}
\end{figure}

\begin{figure}[!t]
\centering
\includegraphics[width=\columnwidth]{figures/node_prototype.jpg}
\caption{Assembled sensor node prototype showing the ESP32 development board, MyoWare~2.0 muscle sensor, TP4056 Li-Po charger, 3.3\,V LDO regulator, decoupling capacitors (C1, C2), and the 500\,mAh Li-Po cell. The green power LED (R1\,=\,470\,$\Omega$) is visible at the top right. The complete node is compact enough to be worn directly over the target muscle belly.}
\label{fig:prototype}
\end{figure}


\subsection{Electrode Configuration and Placement}

Electrodes are mounted on the MyoWare 2.0 snap connectors and positioned according to SENIAM guidelines~\cite{Hermens2000SENIAM}. For the four-channel configuration used here, the target muscles are the rectus femoris, biceps femoris, gastrocnemius medialis, and tibialis anterior---a combination that captures the primary flexor--extensor dynamics of the lower limb during gym-based exercise (squats, lunges, leg press). This muscle selection is consistent with recent eight-channel lower-limb sEMG benchmarks for motion recognition~\cite{mot_semg2025, vijayvargiya2022semg}. The reference electrode for each node is placed over the nearest palpable bony landmark: the lateral femoral condyle for the thigh nodes and the lateral malleolus for the shank nodes. All electrodes are single-use Ag/AgCl disposable pads and are replaced between subjects to prevent cross-contamination. Fig.~\ref{fig:placement} shows the four sensor placement sites on the lower limb, with a close-up of one mounted node in Fig.~\ref{fig:placement_zoom}.

\begin{figure}[!t]
\centering
\includegraphics[width=\columnwidth]{figures/electrode_placement_overview.jpg}
\caption{Placement of all four wireless sensor nodes on the lower limb. From left to right: rectus femoris (anterior thigh), biceps femoris (posterior thigh), tibialis anterior (anterior shank), and gastrocnemius medialis (posterior shank). Each node is secured with an elastic strap; electrode snap connectors are placed directly on the skin according to SENIAM guidelines.}
\label{fig:placement}
\end{figure}

\begin{figure}[!t]
\centering
\includegraphics[width=0.75\columnwidth]{figures/electrode_placement_detail.jpg}
\caption{Close-up of a sensor node mounted over the rectus femoris. The three Ag/AgCl snap electrodes of the MyoWare~2.0 are visible: two differential recording pads aligned along the muscle belly and one passive reference pad over the lateral femoral condyle. The green power LED confirms that the 3.3\,V supply is active.}
\label{fig:placement_zoom}
\end{figure}


\subsection{Signal Acquisition, Transmission, and Host Interface}

Once a 50-sample window is accumulated, the slave firmware packs it into a batch packet and dispatches it to the master's registered MAC address over ESP-NOW. A critical implementation detail concerns the \texttt{t0\_ms} field: rather than deriving this value from each slave's independent \texttt{millis()} counter---which begins at zero at each node's own power-on instant and therefore diverges across devices---the firmware computes
\begin{equation}
  \texttt{t0\_ms} = \texttt{frame\_id\_start} \times 1\,\text{ms},
\end{equation}
where \texttt{frame\_id\_start} is the master's global beacon counter at the first sample of the batch. Because every slave uses the same beacon counter as its time reference, all four channels share a common epoch, making row-wise alignment in the recorded CSV exact without any post-hoc correction.

The master ESP32 receives the batch packets over ESP-NOW and serialises each into a single JSON line forwarded over USB-UART at 921600\,baud:
\begin{verbatim}
{"slave":0,"frame_id_start":1200,
 "t0":1200,"dt_us":1000,"mv":[...50 values...]}
\end{verbatim}
The master requires no WiFi router connection; it runs WiFi in station mode solely to satisfy the ESP-NOW radio driver, and the channel is hard-coded to channel~6 to match all slave nodes. No TCP/IP stack, WebSocket server, or network infrastructure is involved at the host interface.

The host PC reads the serial stream directly through the\,Web\,Serial\,API~\cite{webserial_spec} in the monitoring GUI (described below), which parses the JSON frames in the browser, renders live channel traces, and exports session data to disk in CSV format.

\begin{figure}[!t]
\centering
\includegraphics[width=\columnwidth]{figures/signal_pipeline.png}
\caption{End-to-end signal acquisition pipeline. At each sensor node the MyoWare~2.0 conditions and amplifies the raw sEMG; the ESP32 digitises each beacon-triggered sample on GPIO34 (ADC1\_CH6) at 1\,kSa/s, accumulates 50-point batches, and transmits a $\sim$116-byte ESP-NOW frame to the master. The master forwards the merged stream to the host PC at 921600\,baud over USB-UART, where the browser-based monitoring GUI renders live traces and exports session logs.}
\label{fig:pipeline}
\end{figure}


\subsection{Browser-Based Monitoring GUI}

The host interface is implemented as a single-page web application served from a static hosting platform (Vercel) and accessed in any Chromium-based browser (Google Chrome, Microsoft Edge) that supports the Web Serial API. No native application, driver installation, or Python environment is required on the host machine.

Upon clicking \emph{Connect}, the browser presents the OS-level USB device picker; once the user selects the master ESP32, the dashboard opens a serial port at 921600\,baud and begins parsing the incoming JSON stream entirely in JavaScript. The GUI provides the following functions:

\begin{itemize}
  \item \textit{Real-time waveform display.} Four independent strip charts (one per channel) update at $\sim$30\,frames per second using Chart.js. Each chart displays the most recent 500 samples ($\approx$1\,s of data) with per-channel RMS, peak, peak-to-peak, and mean amplitude readouts updated continuously. An auto-scale toggle fits the Y-axis to the current signal range.

  \item \textit{Software signal conditioning.} A causal biquad IIR filter bank (implemented in the \texttt{emg-engine.js} module) applies a 20\,Hz high-pass, a 450\,Hz low-pass, and a 50\,Hz notch filter to each channel in real time at the correct Nyquist frequency of 1\,kHz~\cite{robertson2004research}. The filter can be toggled on or off at any point; toggling on resets all filter state to prevent transient artefacts.

  \item \textit{Session recording and export.} Participant metadata (ID, sex, age, mass, height) and exercise parameters (type, trial number) are entered via the control panel before each trial. Clicking \emph{Record} captures raw ADC millivolt values together with hardware-derived timestamps for all active channels. On stopping, three CSV formats are available for download: (i)~\emph{Aligned (Filtered)}---a four-column wide-format table in which each row corresponds to a single beacon frame, channels are time-aligned by \texttt{frame\_id}, and all values have been passed through the offline zero-phase filter; (ii)~\emph{Aligned (Raw)}---identical layout without filtering; and (iii)~\emph{Long CSV}---one row per sample across all channels, preserving individual hardware timestamps and sync keys for IEEE-style signal-processing analysis.

  \item \textit{System diagnostics.} A statistics bar displays the cumulative packet count, parse error count, data volume, and effective sample rate (derived from the \texttt{dt\_us} field in each JSON frame rather than browser timing, so all four channels report identically 1\,kHz).

  \item \textit{EMG Biofeedback Game.} A companion page (\emph{Muscle Rush}) uses the same \texttt{emg-engine.js} pipeline to drive a browser-based game: the player controls an in-game character by contracting specific muscles, with difficulty scaling proportional to the channel RMS amplitude. This mode can serve as a motivational tool during data collection sessions or as a standalone biofeedback interface.
\end{itemize}

\begin{figure}[!t]
\centering
\includegraphics[width=\columnwidth]{figures/emg_gui.png}
\caption{Screenshot of the GymEMG-Net browser-based monitoring dashboard. The top panel shows the connection controls, session metadata fields, and recording buttons. The four strip charts display real-time sEMG waveforms for all four channels with per-channel amplitude statistics. The status bar at the bottom reports packet statistics and effective sample rate.}
\label{fig:gui}
\end{figure}



\subsection{Wireless Communication and Node Synchronisation}

ESP-NOW is a proprietary Espressif protocol that encapsulates short data payloads in IEEE\,802.11 vendor-specific action frames and delivers them peer-to-peer---typically at 1\,Mbps in 802.11b compatibility mode---without a TCP/IP stack, association handshake, or routing layer~\cite{espressif_espnow, amponis2022espnow}. The resulting end-to-end latency of 1--10\,ms is roughly two orders of magnitude lower than a Wi-Fi/MQTT path (50--500\,ms)~\cite{eridani2021espnow}. While this latency advantage is less critical for the offline classification pipeline presented in the following section, it would become significant if the system were later adapted for closed-loop real-time biofeedback.

Each slave node is pre-registered to the master's MAC address at firmware initialisation, implementing a one-to-many star configuration consistent with the approach described by Labib~\textit{et~al.}~\cite{labib2021espnow}. Based on published characterisation of comparable ESP-NOW deployments in similar indoor environments, packet loss is expected to remain below 1\% over the operating ranges encountered in a gym setting~\cite{george2023indoor}.

% [AUTHOR NOTE: Consider reporting a measured packet-loss figure from a
% bench test at the actual operating range (e.g., a simple frame counter
% comparison over N sessions). Even a short 5-minute benchtop measurement
// would replace the estimated figure above with a verifiable result.]

\textit{Synchronisation.} Hardware-level inter-node synchronisation is achieved through a master-broadcast beacon mechanism. The master ESP32 emits a 13-byte sync beacon every 1\,ms over the ESP-NOW broadcast address (\texttt{FF:FF:FF:FF:FF:FF}), yielding a system sample rate of 1\,kSa/s per channel. Each beacon contains a monotonically increasing frame counter (\texttt{frame\_id}) and a microsecond-resolution hardware timestamp (\texttt{master\_us}) generated by \texttt{esp\_timer\_get\_time()}. On receipt, each slave node immediately latches its ADC---ensuring that all four channels sample at the same instant, rather than relying on software loop timing. Because the single-sample unicast per beacon was eliminated from the slave firmware (the master forwards batch packets only), no TDMA slot delay is required; batch packets are assembled in a double-buffered ring and dispatched from the Arduino \texttt{loop()} task, fully outside the receive callback. This beacon-driven approach eliminates the inter-channel skew associated with unsynchronised free-running clocks, and the \texttt{frame\_id} field provides the host with an unambiguous sample alignment index without any post-hoc correction. Specifically, the timestamp assigned to each recorded sample is $\texttt{ts}_{i} = (\texttt{frame\_id\_start} + i) \times 1\,\text{ms}$, placing every channel in the same temporal coordinate system and enabling exact row-wise alignment in the exported CSV.

% [AUTHOR NOTE: Consider reporting the measured inter-channel timing jitter
% under this beacon-TDMA scheme, e.g., by logging paired timestamps from
% two nodes simultaneously over a short test session.]


\subsection{Mechanical Enclosure}

Each assembled sensor node is housed in a custom enclosure fabricated using fused deposition modelling (FDM) with polylactic acid (PLA) filament at a layer resolution of 0.2\,mm, providing an acceptable balance between surface finish quality and print time. The enclosure measures 55\,mm\,$\times$\,39\,mm\,$\times$\,22\,mm (length\,$\times$\,width\,$\times$\,height) and consists of two interlocking halves secured by four M2 machine screws, allowing non-destructive access to the internal electronics for firmware updates, battery replacement, or sensor servicing. Cable entry ports with smooth radiused edges protect the sensor cables from chafing and abrasion during wearable use. The enclosure is fitted with an elastic strap that wraps around the limb segment, holding the node firmly in place without restricting joint movement. The assembled enclosure is shown in Fig.~\ref{fig:enclosure}.

\begin{figure}[!t]
\centering
\includegraphics[width=0.75\columnwidth]{figures/node_enclosure.jpg}
\caption{3D-printed PLA enclosure (55\,mm\,$\times$\,39\,mm\,$\times$\,22\,mm) for a single sensor node, shown with the elastic mounting strap. The two interlocking halves are secured by four M2 screws. MyoWare~2.0 electrode snap connectors protrude through openings on the underside for direct skin contact; the green power LED is visible through a small aperture on the top face.}
\label{fig:enclosure}
\end{figure}


\subsection{Power Supply and Safety}

Every sensor node is powered exclusively from its own 3.7\,V / 500\,mAh Li-Po cell, maintaining full galvanic isolation between the subject and any mains-referenced equipment---a fundamental safety requirement for wearable biopotential acquisition systems~\cite{vijayvargiya2022hardware}. The master node communicates with the host PC over USB-UART, but carries no electrode connection to the body, so the four sensor nodes remain electrically isolated from the mains path throughout operation. Ag/AgCl electrodes are used as single-use, disposable components and are replaced between participants.

Based on the quiescent current specifications of the ESP32-WROOM-32 in active transmit mode and the MyoWare 2.0 front-end supply current, the average node current draw is estimated at approximately 85\,mA. Applying a standard 80\% derating factor for Li-Po discharge efficiency, the 500\,mAh cell is expected to provide approximately 4.7\,hours of continuous operation---sufficient for a full training session. Battery voltage is sampled on ADC1\_CH7 once per second and encoded in the status byte appended to every transmitted frame, enabling post-hoc detection of any supply-induced baseline drift.

% [AUTHOR NOTE: The current draw estimate above (85 mA) and the resulting
% 4.7-hour runtime should be verified with a bench measurement (e.g., a USB
% power monitor or inline ammeter) before submission. Report the measured
% draw breakdown: ESP32 in transmit mode alone vs. MyoWare front-end alone.
% This will also allow you to confirm the TP4056 charge current against
% the 4.2 kΩ programming resistor value.]]


\subsection{Design Trade-offs and Limitations}

Distributing acquisition across four independent wireless nodes resolves the cable and motion-artefact problems of centralised architectures, but introduces trade-offs that are worth stating explicitly.

\begin{itemize}

  \item \textit{Sampling rate.} The beacon fires every 1\,ms, yielding a system sample rate of 1\,kSa/s per channel. The MyoWare~2.0 hardware band-pass passes 20--500\,Hz; at $f_s = 1\,\text{kSa/s}$ the Nyquist frequency is exactly 500\,Hz, satisfying the sampling theorem with no margin to spare. A 450\,Hz software low-pass filter applied in the browser provides a 50\,Hz guard band that prevents any spectral leakage from the 450--500\,Hz roll-off region from contaminating the recorded signal. Should a different temporal resolution be required, only \texttt{BEACON\_INTERVAL\_US} in \texttt{ESP\_master.ino} and \texttt{DT\_US} in \texttt{ESP\_slave.ino} need to be adjusted---both the firmware timestamp logic and the browser filter bank adapt automatically via the \texttt{dt\_us} field in each JSON frame.

  \item \textit{Collision avoidance for batch transmissions.} Because single-sample unicast packets per beacon were eliminated in the final firmware architecture, no TDMA slot delay is required. Batch packets (one per 50-sample window) are dispatched from the slave \texttt{loop()} task, which fires every 50\,ms per channel. The ESP-NOW physical layer uses IEEE\,802.11 CSMA-CA for medium access, which handles the low-rate simultaneous transmissions from four nodes with negligible packet loss in a controlled indoor environment~\cite{george2023indoor}. Expanding to a fifth channel would not require any timing changes.

  \item \textit{Battery endurance.} The estimated 4.7-hour runtime covers typical gym sessions but falls short of full-day ambulatory monitoring. Extended deployment would require a higher-capacity cell or a duty-cycling strategy that suspends the radio between transmission windows.

  \item \textit{Software-defined signal conditioning.} Routing the pre-rectification RAW signal to the ADC shifts band-pass filtering, artefact rejection, and feature extraction to software on the host. This preserves spectral information that hardware envelope detection would discard~\cite{phinyomark2018feature}, but increases the processing load on the offline pipeline and requires careful filter design to avoid group-delay artefacts in non-stationary recording conditions~\cite{emg_classification2026}.

  \item \textit{ADC non-linearity.} The ESP32 ADC1 exhibits moderate integral non-linearity, particularly near the supply rails. The 3.3\,V LDO and the 100\,nF bypass capacitor (C2) reduce supply noise, but applications requiring high absolute amplitude accuracy may benefit from an external ADC with better linearity specifications.

\end{itemize}


\subsection{Hardware Summary}

The complete four-node acquisition system plus master gateway was assembled at a total materials cost below \$40\,USD based on component list prices, making it substantially more accessible than commercial multichannel sEMG platforms while remaining fully documented for independent replication. Table~\ref{tab:bom} itemises the principal components for a single sensor node.

\begin{table}[!t]
\renewcommand{\arraystretch}{1.25}
\caption{Bill of Materials for One Sensor Node}
\label{tab:bom}
\centering
\begin{tabular}{lll}
\toprule
\textbf{Component}    & \textbf{Part / Value}            & \textbf{Function}                      \\
\midrule
sEMG front-end        & MyoWare 2.0 (Advancer Tech.)     & Differential amplification, filtering  \\
Microcontroller       & ESP32 DevKit (WROOM-32)          & 12-bit ADC, ESP-NOW slave              \\
Voltage regulator     & 3.3\,V LDO                       & Stable supply for ADC \& front-end     \\
Bulk capacitor (C1)   & 10\,$\mu$F electrolytic          & Low-freq.\ transient suppression       \\
Bypass capacitor (C2) & 100\,nF ceramic                  & HF noise decoupling at ADC supply      \\
Battery               & Li-Po 3.7\,V / 500\,mAh         & Galvanically isolated power supply     \\
Charger module        & TP4056 (R2\,=\,4.2\,k$\Omega$)  & CC/CV Li-Po charging via USB-C         \\
Power switch (S1)     & SPST toggle / slide              & Controls power to ESP32 \& MyoWare     \\
Power indicator LED   & 3\,mm green, 3.3\,V rail         & Visual power-on confirmation           \\
Current resistor (R1) & 470\,$\Omega$                    & LED current limiter                    \\
Electrodes            & Ag/AgCl snap (disposable)        & Skin--sensor interface                 \\
\bottomrule
\end{tabular}
\end{table}